\chapter{Related Work}

This chapter surveys prior work on automated fake news detection and situates this thesis within the existing literature. The review is organized around three core directions relevant to this work: (i) deep learning and LLM-based detection, (ii) domain-aware and 
evidence-grounded classification, and (iii) agentic and multi-agent systems. For each direction, key contributions are summarized and their limitations identified. The chapter concludes by identifying the research gap that this thesis addresses.

\section{Deep Learning and LLM-Based Detection}

Early fake news detection systems relied on manually engineered linguistic and psychological features combined with classical machine learning classifiers such as logistic regression and support vector machines. While effective on early benchmarks, these approaches suffered from limited generalization and vulnerability to sophisticated misinformation strategies. Neural network-based models, including convolutional and recurrent architectures, improved representation learning by capturing contextual and sequential information. The introduction of pre-trained language models such as BERT further advanced detection performance by enabling contextualized text understanding, and hybrid models combining LLM representations with neural classifiers achieved strong results across multiple benchmarks \cite{10.1145/3728725.3728739}.

More recent research explores the direct use of LLMs for misinformation detection. However, studies show that misinformation generated by LLMs is often harder for both humans and automated detectors to identify than human-written misinformation with equivalent semantics, revealing critical weaknesses in existing detection systems \cite{chen2024llmgeneratedmisinformationdetected}. To address data scarcity and class imbalance, some approaches use LLMs to generate synthetic fake news examples and apply reinforcement learning-based sampling strategies, achieving substantial improvements in detection performance \cite{tong-etal-2025-generate}.

Several studies also highlight limits of text-only classification. Diagnostic evaluations on political claim benchmarks suggest a performance ceiling for fine-grained veracity prediction when models rely mostly on lexical or shallow semantic cues \cite{hasan2025generalizationgapspoliticalfake}. Work in applied settings similarly shows that textual signals often dominate prediction, while simple metadata features can still help when language models are not available or computation is limited \cite{savatteri2025signalsreallymattermisinformation}. Together, these findings motivate approaches that incorporate external evidence and focus on robustness, instead of relying only on monolithic content classifiers.

\section{Domain-Aware and Evidence-Grounded Classification} 

Misinformation characteristics vary significantly across domains such as politics, healthcare, and finance, motivating domain-aware and evidence-based detection approaches. Many systems incorporate external knowledge sources, document-level evidence, knowledge graphs, and credibility cues to improve factual verification. Knowledge-enhanced models and graph-based approaches are particularly effective in high-stakes domains by modeling entity--fact relationships and source credibility \cite{10.1145/3728725.3728739}. 

Beyond static knowledge integration, document-grounded factuality checking has emerged as a core component for verifying LLM outputs and detecting misinformation claims. Efficient fact-checkers such as MiniCheck demonstrate that small, specialized sequence-to-sequence models can approximate large-model verification quality when trained on high-quality synthetic errors, reducing inference cost and latency while remaining evidence-grounded \cite{tang-etal-2024-minicheck}. In parallel, modular toolkits and unifying frameworks (e.g., OpenFactCheck) aim to standardize the construction and evaluation of fact-checking pipelines and LLM factuality assessment, improving comparability across methods \cite{wang2025openfactcheckbuildingbenchmarkingcustomized}. Survey work further consolidates this space by categorizing how generative LLMs are used across the fact-checking pipeline, including evidence retrieval, claim verification, and synthetic data generation \cite{vykopal2024generativelargelanguagemodels}. 

Domain-specific pipelines also increasingly combine retrieval quality signals with selective reasoning. In health misinformation, two-stage systems first estimate agreement or stance across retrieved evidence and only trigger more expensive multi-agent deliberation when evidence is conflicting or uncertain, improving efficiency while maintaining accuracy \cite{chen2025exploringhealthmisinformationdetection}. More generally, confidence-aware routing approaches argue that retrieval should be invoked adaptively (rather than indiscriminately), using signals such as probabilistic certainty and reasoning consistency to decide between internal knowledge, targeted retrieval, or deeper search \cite{wang2026factcheckinglargelanguagemodels}. 

\section{Multi-Agent, Agentic, and Ensemble Approaches} 

To overcome limitations of monolithic detection systems, recent work has explored agentic and multi-agent frameworks for misinformation detection and fact-checking. Single-agent pipelines such as FactAgent simulate expert fact-checking workflows by decomposing verification into multiple structured reasoning steps, including internal knowledge checks and external evidence retrieval. These systems improve transparency and performance without requiring supervised training \cite{li2024largelanguagemodelagent}. 

Multi-agent systems extend this paradigm by assigning specialized roles to multiple agents that collaborate or compete to assess a claim's veracity. Debate-based frameworks such as TruEDebate frame fake news detection as a structured argumentative process, where agents generate opposing arguments and synthesis agents aggregate them into a final decision, improving robustness and interpretability \cite{10.1145/3726302.3730092}. Other multi-agent frameworks adopt modular architectures spanning the misinformation lifecycle, including detection, correction, source verification, and provenance tracking \cite{gautam2025multiagentsystemsmisinformationlifecycle}. Logical and causal reasoning can also be explicitly introduced via specialized agent roles (e.g., counterfactual or causal evaluators) to detect inconsistencies and improve reasoning transparency \cite{10.1145/3696410.3714748}. Relatedly, distributed decision formulations study how to aggregate labels from multiple imperfect agents by learning agent reliabilities over time, connecting misinformation verification to online learning and crowdsourcing-style aggregation \cite{verma2025multiagentfactchecking}. 

A recurring theme in agentic systems is the tight coupling between debate and retrieval. Tool-augmented debate frameworks equip different agents with heterogeneous tools (e.g., web search vs.\ curated RAG corpora) and iteratively refine queries as the debate evolves, reducing hallucinations and improving grounding \cite{jeong2026toolmadmultiagentdebateframework}. Similar ideas appear in zero-shot fake news detection under knowledge cutoff constraints: retrieval is driven by entity-based query generation, followed by multi-LLM role-based interaction and adjudication to produce interpretable verdicts \cite{wu2025zofiazeroshotfakenews}. For multimodal and real-world settings, evidence-grounded, multi-persona agentic frameworks further combine web evidence, reverse image search, and knowledge graph paths, often paired with a lightweight classifier to mitigate LLM uncertainty \cite{bukke2025agenticmultipersonaframeworkevidenceaware}. 

\section{Robustness, Adversarial Misinformation, and Style Invariance} 

As adversaries increasingly mimic credible writing style and exploit generative models, robustness has become a central concern. Empirical evidence suggests LLM-generated misinformation can evade both human judgment and automated detectors more effectively than comparable human-written misinformation \cite{chen2024llmgeneratedmisinformationdetected}. To reduce reliance on style cues, several methods shift toward content-grounded or event-centric representations. For example, FACTGUARD extracts event-centric content using LLMs and applies a reliability mechanism to down-weight contradictory or unclear advice; it also includes a distilled variant for resource-constrained deployment \cite{he2025factguardeventcentriccommonsenseguidedfake}. Complementarily, adversarial co-evolution frameworks such as RADAR train a generator to produce harder fake variants via factual perturbations while improving a retrieval-augmented detector, using natural-language critiques (verbal adversarial feedback) to strengthen robustness beyond scalar reward signals \cite{ma2026radarretrievalaugmenteddetectoradversarial}. 

\section{Multimodality, Propagation, Early Detection, and Fine-Grained Benchmarks} 

Real-world misinformation is often multimodal and time-sensitive. Multimodal systems integrate text with images and other contextual signals to detect cross-modal inconsistencies and improve robustness \cite{10.1145/3728725.3728739}. Early detection is especially challenging because propagation signals may not yet be observable; recent agent-driven approaches address this by simulating early diffusion patterns (`virtual propagation'') with role-differentiated LLM agents and fusing the simulated social evidence with content signals, improving performance in cold-start stages \cite{gu2026aheadspreadagentdrivenvirtual}. 

Finally, newer benchmarks highlight that binary document-level labels can hide fine-grained falsity patterns. Span-level datasets such as MisSpans annotate which exact spans are false within sentences and categorize misinformation types, enabling more actionable explanations and revealing that fine-grained localization remains difficult for current LLMs even when coarse classification appears strong \cite{liu2026misspansfinegrainedfalsespan}. Related benchmark work also targets ambiguity in intent, such as distinguishing satire from fake news, and shows that lightweight transformer models can be competitive while emphasizing calibration and statistical testing \cite{chhetri2025wisewebinformationsatire}. 

\section{Evaluation Frameworks and End-to-End Fact-Checking Assessment} 

Beyond model architectures, evaluation methodology strongly shapes conclusions. Arena-style evaluation frameworks benchmark end-to-end fact-checking pipelines (claim extraction/decomposition, tool-augmented retrieval, and justification-driven verdict prediction) and show that strong claim-level accuracy may not translate into full-pipeline competence \cite{lin2026comprehensivestagewisebenchmarkinglarge}. Similarly, adaptive evaluation approaches generate model-centric test cases iteratively to probe reasoning quality and justification faithfulness beyond static benchmarks \cite{lin-etal-2025-fact}. In applied political fact-checking, large-scale studies indicate that web-search or reasoning modes yield only modest gains unless retrieval is curated and claim-specific, underscoring the importance of retrieval quality and evaluation design \cite{deverna2025largelanguagemodelsrequire}. 

Overall, recent trends move from monolithic text classifiers toward evidence-grounded and agentic systems that integrate retrieval, 
specialized reasoning roles, and calibrated decision aggregation. While these approaches improve interpretability and robustness, they also introduce new challenges in coordination complexity, cost, and trustworthy reporting of evidence and system success 
\cite{DBIS2025,chen2025evidenceboundautonomousresearchevibound}. Despite this progress, several gaps relevant to short political claim verification remain, which are addressed in the following section.

\section{Research Gap and Positioning of This Thesis}

Despite substantial progress, several important gaps remain in the existing literature. First, most LLM-based detection systems are 
evaluated on general or long-form content and do not specifically address the challenges of short political claims, where the veracity signal is weak and domain-specific knowledge is critical. Second, while domain-aware approaches have been explored, few studies systematically investigate how different levels of domain granularity affect both in-domain performance and cross-domain generalization on a political fact-checking benchmark such as LIAR. Third, existing multi-agent systems for fact-checking predominantly rely on external retrieval or web search, whereas the feasibility and effectiveness of retrieval-free, domain-specialized multi-agent architectures remains underexplored.

This thesis addresses these gaps by systematically investigating binary claim verification on the LIAR benchmark using retrieval-free LLM-based approaches. Concretely, it evaluates how domain knowledge incorporated through prompting, fine-tuning, and agent specialization affects classification accuracy and cross-domain generalization (RQ1a, RQ1b), and whether multi-agent architectures with domain routing and aggregation mechanisms outperform single-model baselines in this setting (RQ2). Unlike most prior work, this thesis focuses on the interaction between domain specialization and multi-agent coordination without relying on external evidence retrieval, making it directly comparable to prior prompt-based and fine-tuning baselines on the same benchmark.